\section{Methodology: Riemannian Isometry-respecting Manifold Operations}
\label{sec:method}

We present the mathematical framework of \textbf{RIMO} (\textbf{R}iemannian \textbf{I}sometry-respecting \textbf{M}anifold \textbf{O}perations) to merge models on the manifold of the orthogonal group $\mathrm{O}(d)$ and its associated Lie algebra $\mathfrak{so}(d)$. The high-level flowchart of this pipeline is illustrated in Figure~\ref{fig:pipeline_schematic}.

\begin{figure}[ht]
\centering
\begin{tikzpicture}[
  node distance=0.35cm and 0.5cm,
  flat/.style={draw=gray!80, rectangle, rounded corners, fill=gray!20, thick, text width=6.8cm, align=center, font=\small, minimum height=0.6cm},
  group/.style={draw=blue!80, rectangle, rounded corners, fill=blue!15, thick, text width=6.8cm, align=center, font=\small, minimum height=0.6cm},
  algebra/.style={draw=red!80, rectangle, rounded corners, fill=red!15, thick, text width=6.8cm, align=center, font=\small, minimum height=0.6cm},
  arrow/.style={thick, ->, >=stealth}
]
  \node (input) [flat] {\textbf{Input Models} \\ $W_0$ (Base) and $\{W_k\}_{k=0}^{N-1}$ (Task Experts)};
  
  \node (decouple) [group, below=of input] {\textbf{1. Orthogonal Procrustes Decoupling} \\ $W_k \approx W_0 R_k + \rho_k \quad \Longrightarrow \quad R_k \in \mathrm{O}(d), \ \rho_k \in \mathbb{R}^{d \times d}$};
  
  \node (map_tangent) [algebra, below=of decouple] {\textbf{2. Manifold Tangent Projection} \\ $Q_k = (R_k - I_d)(R_k + I_d)^{-1} \in \mathfrak{so}(d)$ \quad (Inv. Cayley)};
  
  \node (agg_spec) [algebra, below=of map_tangent] {\textbf{3. Tangent Aggregation \& Spectral Operations} \\ $Q_{com} = \text{Scale}\left(\sum Q_k\right)$ \\ $\hat{Q}_{com} = \text{Rank-Pruning}(Q_{com})$};
  
  \node (reconstruct) [group, below=of agg_spec] {\textbf{4. Manifold Retraction \& Residual Merging} \\ $R_{merged} = (I + \hat{Q}_{com})(I - \hat{Q}_{com})^{-1} \in \mathrm{O}(d)$ \quad (Cayley) \\ $\rho_{merged} = \frac{1}{N}\sum \rho_k$};
  
  \node (output) [flat, draw=green!80, fill=green!20, below=of reconstruct] {\textbf{Final Merged Model} \\ $W_{final} = W_0 R_{merged} + \rho_{merged}$};

  \draw [arrow] (input) -- (decouple);
  \draw [arrow] (decouple) -- (map_tangent);
  \draw [arrow] (map_tangent) -- (agg_spec);
  \draw [arrow] (agg_spec) -- (reconstruct);
  \draw [arrow] (reconstruct) -- (output);
\end{tikzpicture}
\caption{\textbf{Schematic of the RIMO Model Merging Pipeline.} Standard linear task-specific weights are decoupled into geometric (orthogonal) and linear (residual) components, mapped and operated on inside the Lie algebra tangent space, and retracted back to reconstruct the final model. Color coding: light gray represents standard Euclidean parameter space, light blue represents the orthogonal Lie Group manifold, and light red represents the Lie algebra tangent space.}
\label{fig:pipeline_schematic}
\vskip -0.1in
\end{figure}

Let $W_0 \in \mathbb{R}^{d \times d}$ be the pre-trained base weight matrix, and let $\{W_k\}_{k=0}^{N-1}$ be the $N$ fine-tuned expert weights. For rectangular weights, we apply block-diagonal operations on square sub-matrices of size $b \times b$ \cite{oft2023, yang2026orthomerge}. Although block-diagonal partitioning restricts the orthogonal degrees of freedom to localized coordinate compartments (bypassing global layer rotations), it introduces zero coordinate clipping or boundary distortions because the localized Cayley transforms operate independently on separate coordinate sub-blocks, preserving their algebraic orthogonality and block-wise boundaries exactly (further analyzed in Appendix \ref{subsec:rectangular_boundaries}).

The complete end-to-end execution of RIMO and RIMO-Pruned is formalized in Algorithm~\ref{alg:rimo_pruned_algorithm}.

\begin{algorithm}[tb]
   \caption{RIMO and RIMO-Pruned Model Merging}
   \label{alg:rimo_pruned_algorithm}
\begin{algorithmic}[1]
   \STATE {\bfseries Input:} Pre-trained base weights $W_0 \in \mathbb{R}^{d \times d}$, fine-tuned expert weights $\{W_k\}_{k=0}^{N-1}$, keep ratio $r \in (0, 1]$, and residual scale $s \in [0, 1]$.
   \STATE {\bfseries Output:} Final merged weights $W_{\text{final}} \in \mathbb{R}^{d \times d}$.
   \STATE Initialize $R_{\text{merged}} \leftarrow I_d$ and $\rho_{\text{merged}} \leftarrow 0$.
   \FOR{each expert $k = 0, \dots, N-1$}
     \STATE Compute right Orthogonal Procrustes decoupling:
     \STATE \quad $U_k \Sigma_k V_k^T = \mathrm{SVD}(W_0^T W_k)$
     \STATE \quad $R_k = U_k V_k^T \in \mathrm{O}(d)$
     \STATE \quad $\rho_k = W_k - W_0 R_k$
     \STATE Map rotation $R_k$ to Lie algebra tangent space $\mathfrak{so}(d)$ via Inverse Cayley:
     \STATE \quad $Q_k = (R_k - I_d)(R_k + I_d)^{-1}$
   \ENDFOR
   \STATE Compute magnitude-corrected Lie algebra average $Q_{\text{com}}$:
   \STATE \quad $Q_{\text{avg}} = \frac{1}{N} \sum_{k=0}^{N-1} Q_k$
   \STATE \quad $c = \sum_{k=0}^{N-1} \|Q_k\|_F / \|\sum_{k=0}^{N-1} Q_k\|_F$
   \STATE \quad $Q_{\text{com}} = c \cdot Q_{\text{avg}}$
   \IF{using RIMO (Isotropic Spectral Balancing)}
     \STATE Perform SVD: $U_{\text{com}} \Sigma_{\text{com}} V_{\text{com}}^T = \mathrm{SVD}(Q_{\text{com}})$ where $\Sigma_{\text{com}} = \text{diag}(\sigma_1, \dots, \sigma_d)$
     \STATE Compute mean singular value: $\bar{\sigma} = \frac{1}{d} \sum_{i=1}^d \sigma_i$
     \STATE Balance spectrum towards mean: $\hat{\Sigma}_{\text{com}} = \bar{\sigma} I_d + (\Sigma_{\text{com}} - \bar{\sigma} I_d) / \sqrt{N}$
     \STATE Reconstruct: $Q'_{\text{com}} = U_{\text{com}} \hat{\Sigma}_{\text{com}} V_{\text{com}}^T$
     \STATE Skew-symmetric project: $\hat{Q}_{\text{com}} = \frac{1}{2}(Q'_{\text{com}} - (Q'_{\text{com}})^T)$
   \ELSIF{using RIMO-Pruned (Rank-Preserving Spectral Pruning)}
     \STATE Perform SVD: $U_{\text{com}} \Sigma_{\text{com}} V_{\text{com}}^T = \mathrm{SVD}(Q_{\text{com}})$ where $\Sigma_{\text{com}} = \text{diag}(\sigma_1, \dots, \sigma_d)$
     \STATE Select top $\lfloor r \cdot \text{rank}(Q_{\text{com}}) \rfloor$ singular value pairs in $\Sigma_{\text{com}}$, zeroing the rest to form $\hat{\Sigma}_{\text{com}}$
     \STATE Reconstruct: $Q'_{\text{com}} = U_{\text{com}} \hat{\Sigma}_{\text{com}} V_{\text{com}}^T$
     \STATE Skew-symmetric project: $\hat{Q}_{\text{com}} = \frac{1}{2}(Q'_{\text{com}} - (Q'_{\text{com}})^T)$
   \ELSE
     \STATE Keep un-modified generator: $\hat{Q}_{\text{com}} = Q_{\text{com}}$
   \ENDIF
   \STATE Map generator $\hat{Q}_{\text{com}}$ back to Lie Group $\mathrm{O}(d)$ via forward Cayley transform:
   \STATE \quad $R_{\text{merged}} = (I_d + \hat{Q}_{\text{com}})(I_d - \hat{Q}_{\text{com}})^{-1}$
   \STATE Merge residual components: $\rho_{\text{merged}} = \frac{s}{N} \sum_{k=0}^{N-1} \rho_k$
   \STATE Reconstruct final weights: $W_{\text{final}} = W_0 R_{\text{merged}} + \rho_{\text{merged}}$
\end{algorithmic}
\end{algorithm}

\subsection{Phase 1: Orthogonal Extraction via Procrustes Decomposition}
Since standard fine-tuning is conducted in Euclidean space, expert weights $\{W_k\}$ are generally non-orthogonal. We decouple them into an orthogonal rotation $R_k \in \mathrm{O}(d)$ and a residual component $\rho_k \in \mathbb{R}^{d \times d}$ using the right-multiplication \emph{Orthogonal Procrustes problem}:
\begin{equation}
\label{eq:procrustes}
R_k = \arg\min_{R} \|W_k - W_0 R\|_F^2 \quad \text{s.t.} \quad R^T R = I_d
\end{equation}
Expanding the objective, we obtain:
\begin{equation*}
\|W_k - W_0 R\|_F^2 = \|W_k\|_F^2 + \|W_0\|_F^2 - 2 \mathrm{Tr}(W_k^T W_0 R)
\end{equation*}
To minimize this, we maximize $\mathrm{Tr}(W_0^T W_k R^T)$. Let the SVD of the cross-covariance be: $U_k \Sigma_k V_k^T = \mathrm{SVD}(W_0^T W_k)$. The closed-form solution to Eq.~\ref{eq:procrustes} is $R_k = U_k V_k^T$. The linear displacement not captured by pure rotation is stored as the residual: $\rho_k = W_k - W_0 R_k$.

---

\subsection{Phase 2: Projecting to the Tangent Space $\mathfrak{so}(d)$}
We map the rotations $\{R_k\}$ to the tangent space at the identity, i.e., the Lie algebra $\mathfrak{so}(d)$ of skew-symmetric matrices, using the \emph{inverse Cayley transform}:
\begin{equation}
\label{eq:inv_cayley}
Q_k = (R_k - I_d)(R_k + I_d)^{-1}
\end{equation}

\begin{proposition}[Algebraic Closure of the Inverse Cayley Transform]
For any orthogonal matrix $R \in \mathrm{O}(d)$ such that $\det(R + I_d) \ne 0$, the matrix $Q = (R - I_d)(R + I_d)^{-1}$ is skew-symmetric ($Q^T = -Q$)\footnote{Mathematically, the inverse Cayley transform is undefined when $R$ has an eigenvalue of $-1$ (i.e., when $\det(R + I_d) = 0$). In practice, numerical stability is easily safeguarded by adding a tiny diagonal perturbation $\epsilon I_d$ (with $\epsilon = 10^{-6}$) to the matrix prior to inversion if any eigenvalues lie near $-1$, or by utilizing the matrix logarithm map (which is globally defined on $\mathrm{SO}(d)$ but is significantly more computationally expensive).}.
\end{proposition}

The proof of Proposition 3.1 is provided in Appendix~\ref{proof:proposition_closure}.

---

\subsection{Phase 3: Magnitude-Corrected Aggregation}
The tangent space $\mathfrak{so}(d)$ is a vector space, allowing linear combination. To prevent destructive interference and preserve rotational magnitude, we apply a magnitude correction:
\begin{equation}
Q_{com} = c \cdot \left( \frac{1}{N} \sum_{k=0}^{N-1} Q_k \right)
\end{equation}
where the scalar correction coefficient is defined inline as $c = \sum_{k=0}^{N-1} \|Q_k\|_F / \left\| \sum_{k=0}^{N-1} Q_k \right\|_F$.

---

\subsection{Phase 4: Isotropic Spectral Balancing in $\mathfrak{so}(d)$}
To prevent single-task dominance, we introduce spectral balancing directly within the Lie algebra $\mathfrak{so}(d)$. We compute the SVD of the aggregated generator $Q_{com} = U_{com} \Sigma_{com} V_{com}^T$, where $\Sigma_{com} = \mathrm{diag}(\sigma_1, \dots, \sigma_d)$ contains the singular values. The mean singular value is computed as the arithmetic average:
\begin{equation}
\label{eq:mean_sigma}
\bar{\sigma} = \frac{1}{d} \sum_{i=1}^d \sigma_i
\end{equation}
We explicitly clarify that this sum in Eq.~\ref{eq:mean_sigma} is taken over all $d$ singular values (including both zero singular values from inactive dimensions and both duplicate entries in active ones), rather than only the distinct or non-zero ones. Because any real skew-symmetric matrix $Q_{com}$ has singular values that naturally appear in duplicate pairs, this formulation ensures that the mean $\bar{\sigma}$ properly acts as a global scaling factor across the entire representational space. This choice is vital: using the full diagonal forces the target isotropic level $\bar{\sigma}$ to scale down proportionally with the number of inactive dimensions, though as we show in Section 4.5, this also causes the noise power under Cayley mapping to scale quadratically with representation size.

We interpolate the singular spectrum towards uniform isotropy with hyperparameter $t \ge 1$:
\begin{equation}
\label{eq:spectral_interpolate}
\hat{\Sigma}_{com} = \bar{\sigma} I_d + (\Sigma_{com} - \bar{\sigma} I_d) \times \frac{1}{\sqrt{t}}
\end{equation}
and reconstruct the balanced matrix $Q'_{com} = U_{com} \hat{\Sigma}_{com} V_{com}^T$. To ensure strict skew-symmetry under floating-point precision, we project it back onto $\mathfrak{so}(d)$:
\begin{equation}
\label{eq:projection}
\hat{Q}_{com} = \frac{1}{2}(Q'_{com} - (Q'_{com})^T)
\end{equation}

\subsubsection{The Kernel Distortion Theorem}
When combining SVD with skew-symmetric projection on Lie algebras, a major mathematical challenge arises in the multi-dimensional null space of the generator. We formalize this via the \emph{Kernel Distortion Theorem}, which demonstrates how standard numerical SVD solver implementations inject destructive skew-symmetric noise when attempting to balance zero singular values.

\begin{theorem}[Kernel Distortion Theorem]
Let $Q_{com} \in \mathfrak{so}(d)$ with a nullity $m = \dim(\ker(Q_{com})) \ge 2$, and let $U \Sigma V^T = \mathrm{SVD}(Q_{com})$. Let $\mathcal{K}$ be the set of indices corresponding to zero singular values (the kernel), and let $Q'_{kernel} = \sum_{i \in \mathcal{K}} \hat{\sigma} U_i V_i^T$ be the reconstructed kernel component inflated to a uniform balanced strength $\hat{\sigma} > 0$. Under the skew-symmetric projection operator $\mathcal{P}(M) = \frac{1}{2}(M - M^T)$, the projected kernel component is given by:
\begin{equation}
\mathcal{P}(Q'_{kernel}) = \frac{\hat{\sigma}}{2} V_{kernel} (P - P^T) V_{kernel}^T
\end{equation}
where $V_{kernel} \in \mathbb{R}^{d \times m}$ is the orthonormal basis of the kernel and $P \in \mathrm{O}(m)$ is an arbitrary orthogonal gauge transformation. Thus, the projection is zero if and only if $P$ is symmetric.
\end{theorem}

The complete proof of Theorem 3.2 is provided in Appendix~\ref{proof:theorem_distortion}. This theorem establishes that standard numerical SVD solvers introduce a non-symmetric gauge transformation $P$ in the multi-dimensional kernel, causing the projection to inject non-zero skew-symmetric noise into the kernel of the generator. This mathematically explains how SVD-based spectral balancing in Lie algebras is fundamentally destructive.

\subsubsection{The Spectrum Distortion Theorem}
Beyond the null space, the projection operator $\mathcal{P}$ also introduces systematic spectral distortion in the active (non-zero) subspace when singular values are modified non-uniformly. We formalize this in the following theorem:

\begin{theorem}[Spectrum Distortion Theorem]
Let $Q \in \mathfrak{so}(d)$ with SVD $Q = U \Sigma V^T$, and let $R = U^T V \in \mathrm{O}(d)$ be the orthogonal matrix relating the left and right singular bases. Let $\hat{\Sigma}$ be a diagonal matrix of modified singular values, and let $Q' = U \hat{\Sigma} V^T$. The skew-symmetric projection of $Q'$ is $\hat{Q} = \mathcal{P}(Q') = \frac{1}{2}(Q' - Q'^T)$.
Then:
\begin{enumerate}
    \item The projected matrix $\hat{Q}$ is skew-symmetric, and its singular values match the target $\hat{\Sigma}$ if and only if the relation $R \hat{\Sigma} = -\hat{\Sigma} R^T$ is satisfied.
    \item If the modification is non-uniform (such as the isotropic interpolation in Eq.~\ref{eq:spectral_interpolate} with $t > 1.0$), then $R \hat{\Sigma} \ne -\hat{\Sigma} R^T$ for any active dimensions with unequal singular values. Consequently, the projection $\mathcal{P}$ distorts the singular value spectrum, pulling the actual singular values of $\hat{Q}$ away from the target $\hat{\Sigma}$.
\end{enumerate}
\end{theorem}

The proof of Theorem 3.3 is provided in Appendix~\ref{proof:theorem_distortion_active}. This theorem reveals a fundamental mathematical inconsistency in any manifold merging approach that attempts to modify Lie algebra singular values via standard SVD: because SVD-based modifications are blind to the underlying skew-symmetry relation $R \Sigma = -\Sigma R^T$, the subsequent projection step required to restore skew-symmetry inevitably distorts the spectrum, preventing the model from achieving its target isotropic distribution. This provides a rigorous mathematical explanation of why standard SVD is structurally incompatible with spectral operations in Lie algebras.

\subsubsection{Mitigation: Symmetry-Preserving Decompositions and Rank-Preserving Operations}
To overcome this SVD-induced kernel distortion and mitigate the spectral balancing pitfall, we propose two distinct architectural pathways for manifold model merging:

\textbf{(1) Symmetry-Preserving Spectral Decomposition:} Instead of standard SVD, we propose performing decompositions that respect the algebraic structure of the Lie algebra. Specifically, any skew-symmetric matrix $Q_{com} \in \mathfrak{so}(d)$ can be decomposed using the \emph{real Schur decomposition}:
\begin{equation}
Q_{com} = X T X^T
\end{equation}
where $X \in \mathrm{O}(d)$ is a strictly orthogonal matrix, and $T \in \mathbb{R}^{d \times d}$ is a block-diagonal matrix of the form:
\begin{equation}
T = \mathrm{diag}(T_1, T_2, \dots, T_m, 0, \dots, 0)
\end{equation}
where each active block $T_j \in \mathbb{R}^{2 \times 2}$ is a skew-symmetric matrix:
\begin{equation}
T_j = \begin{pmatrix} 0 & \lambda_j \\ -\lambda_j & 0 \end{pmatrix}
\end{equation}
and the eigenvalues of $Q_{com}$ are conjugate pairs $\pm i \lambda_j$ (the singular values are exactly duplicate pairs $\{\lambda_j, \lambda_j\}$). To perform spectral balancing or pruning, we can modify the block parameters directly:
\begin{equation}
T_j' = \begin{pmatrix} 0 & \hat{\lambda}_j \\ -\hat{\lambda}_j & 0 \end{pmatrix}
\end{equation}
to obtain the modified block matrix $T'$. Because $T'$ is block-diagonal with skew-symmetric blocks, it is strictly skew-symmetric ($T'^T = -T'$). The reconstructed matrix:
\begin{equation}
\hat{Q}_{com} = X T' X^T
\end{equation}
is mathematically guaranteed to be strictly skew-symmetric ($\hat{Q}_{com}^T = -\hat{Q}_{com}$), and its singular values are exactly the modified target magnitudes $|\hat{\lambda}_j|$. Crucially, this decomposition completely eliminates the need for the post-hoc projection step $\mathcal{P}$, thus preserving both the modified spectrum and the Lie algebra structure by design. Alternatively, complex Hermitian eigendecomposition of $i Q_{com} = W \Lambda W^H$ can be used, which enforces a symmetric gauge choice $U_{kernel} = V_{kernel}$ up to sign ($P = I$), avoiding kernel distortion. Naive SVD solvers in standard deep learning libraries (such as `torch.linalg.svd`), however, are unaware of these geometric structures, necessitating the use of our rank-preserving spectral pruning as a stable Euclidean-compatible proxy.

\textbf{(2) Rank-Preserving Spectral Pruning:} Since inflating small singular values in inactive dimensions injects catastrophic rotational noise, an alternative is to perform \emph{spectral pruning} (or truncation) instead of isotropic smoothing. By setting small singular values $\sigma_i < \epsilon$ to exactly zero, we preserve the low-rank structure of the generator ($Q'_{kernel} = 0$, and thus $\mathcal{P}(Q'_{kernel}) = 0$ is trivially satisfied). This keeps inactive dimensions at exactly zero, preventing the injection of spurious rotations and ensuring representational stability while still smoothing the active singular spectrum.

---

\subsection{Phase 5: Manifold Reconstruction and Residual Merging}
We map the balanced generator $\hat{Q}_{com}$ back to $\mathrm{O}(d)$ via the forward Cayley transform: $R_{merged} = (I_d + \hat{Q}_{com})(I_d - \hat{Q}_{com})^{-1}$. The linear residuals $\{\rho_k\}$ are combined via Euclidean arithmetic: $\rho_{merged} = \frac{1}{N} \sum_{k=0}^{N-1} \rho_k$. Finally, the merged weight matrix $W_{final}$ is reconstructed as:
\begin{equation}
W_{final} = W_0 R_{merged} + \rho_{merged}
\end{equation}
This completes the RIMO merging process.
